%% ---------------------------------------------------------------------------
%% 2. Discrete Galerkin meshfree formulation for Kirchhoff--Love plate
%% ---------------------------------------------------------------------------
\section{Discrete Galerkin meshfree formulation for Kirchhoff--Love plate}\label{sec:formulation}

\subsection{Galerkin weak form of Kirchhoff--Love plate}

In accordance with the Kirchhoff--Love plate theory, the equilibrium equation and boundary conditions of an elastic body $\Omega$ with boundary $\Gamma$ can be stated as follows:
\begin{equation}\label{eq:1}
\begin{cases}
-M_{ij,ij}=q & \text{in }\Omega\\[2pt]
w=\bar{w} & \text{on }\Gamma^{w}\\[2pt]
w_{,n_i}=\bar{\theta}_{n} & \text{on }\Gamma^{\theta}\\[2pt]
M_{ij,j}n_i+M_{ns,s}=\bar{V} & \text{on }\Gamma^{V}\\[2pt]
M_{ij}n_in_j=\bar{M}_{n} & \text{on }\Gamma^{M}\\[2pt]
M_{ns}=\bar{F}_{c} & \text{at }\Gamma^{c}
\end{cases}
\end{equation}
where $M_{ij}=-C_{ijkl}w_{,kl}$ and $q$ represent the moment and loading along the plate.
$\Gamma^{w}$ and $\Gamma^{\theta}$ denote the displacement and slope boundaries, $\bar{w}$ and $\bar{\theta}_{n}$ are the corresponding given values.
$\bm{n}=\{n_i\}$ and $\bm{s}=\{s_i\}$ represent the boundary outward normal and tangent.
$\bar{V}$, $\bar{M}_{n}$ and $\bar{M}_{ns}=M_{jk}s_jn_k$ denote the equivalent shear force, moment and torque on the natural boundaries $\Gamma^{V}$ and $\Gamma^{M}$, respectively.
In addition, $\bar{F}_{c}$ is the concentrated force at the corners of boundary curve $\Gamma^{c}$.

The weak form of Kirchhoff--Love plate can be established as follows:
\begin{equation}\label{eq:2}
\int_{\Omega}\delta w\left(M_{ij,ij}+q\right)\dd\Omega=0
\end{equation}
where $\delta w$ denotes an arbitrary test function or virtual displacement.
By invoking the procedure of integration by parts, Eq. \eqref{eq:2} becomes:
\begin{equation}\label{eq:3}
\int_{\Omega}\delta w_{,ij}M_{ij}\dd\Omega
=-\int_{\Gamma}\delta wM_{ij,j}n_i\dd\Gamma
-\int_{\Omega}\delta wq\dd\Omega
+\int_{\Gamma}\delta w_{,j}M_{ij}n_i\dd\Gamma
\end{equation}
To achieve independent and meaningful boundary conditions, the last term in Eq. \eqref{eq:3} is further divided into two parts:
\begin{equation}\label{eq:4}
\begin{aligned}
\int_{\Gamma}\delta w_{,i}M_{ij}n_j\dd\Gamma
&=\int_{\Gamma}\delta w_{,i}\underbrace{\delta_{ij}}_{=\,n_in_j+s_is_j}M_{jk}n_k\dd\Gamma\\[2pt]
&=\int_{\Gamma}\delta w_{,i}\underbrace{n_i}_{=\,\delta w_{,n}}n_jM_{jk}n_k\dd\Gamma
+\int_{\Gamma}\delta w_{,i}\underbrace{s_i}_{=\,\delta w_{,s}}s_jM_{jk}n_k\dd\Gamma\\[2pt]
&=\int_{\Gamma}\delta w_{,n}\bar{M}_{n}\dd\Gamma
-\int_{\Gamma}\delta w\left(s_jM_{jk}n_k\right)_{,s}\dd\Gamma
+\left[\delta ws_jM_{jk}n_k\right]_{\Gamma}\\[2pt]
&=\int_{\Gamma}\delta w_{,n}\bar{M}_{n}\dd\Gamma
-\int_{\Gamma}\delta wM_{ns,s}\dd\Gamma
+\left[\delta wM_{ns}\right]_{\Gamma}
\end{aligned}
\end{equation}
Meanwhile, with the smooth boundary surface assumption, the second term in the right of Eq. \eqref{eq:4} can be eliminated:
\begin{equation}\label{eq:5}
\left[\delta wM_{ns}\right]_{\Gamma}
=\sum_{m=1}^{n_{c}}\left\{\left[\delta wM_{ns}\right](c_{m+1})-\left[\delta wM_{ns}\right](c_{m})\right\}=0
\end{equation}
where $n_{c}$ represents the number of convex points in closed boundary curve, where the corner index $c_{m+1}$ is identified with corner index $c_{m}$.

In order to obtain the general form without assumption, here we still adopt Eq. \eqref{eq:3} to further deduce the weak form of Kirchhoff--Love plate theory.

Further let $\delta w$ and $\delta w_{,n}$ vanish on $\Gamma^{w}$ and $\Gamma^{\theta}$, the following final weak form is rationally attained from Eqs. \eqref{eq:3} and \eqref{eq:5}:
\begin{equation}\label{eq:6}
\int_{\Omega}\delta w_{,ij}M_{ij}\dd\Omega
=-\int_{\Gamma^{V}}\delta w\bar{V}\dd\Gamma
-\int_{\Omega}\delta wq\dd\Omega
+\int_{\Gamma^{M}}\delta w_{,n}\bar{M}_{n}\dd\Gamma
+\left[\delta w\bar{F}_{c}\right]_{\Gamma^{c}}
\end{equation}
where $\bar{V}=\bar{M}_{ns,s}+M_{ij,j}n_i$.

\subsection{MLS/RK meshfree approximation}\label{sec:mls}

In meshfree approximation, a group of nodes $\{x_I\}_{I=1}^{NP}$ is adopted to discretize the problem domain $\Omega$ and corresponding boundary $\Gamma$, where $NP$ stands for the total number of meshfree nodes.
The meshfree approximation of deflection $w(x)$ $\approx$ $w^{h}(x)$ is given by:
\begin{equation}\label{eq:7}
w^{h}(x)=\sum_{I=1}^{NP}\Psi_I(x)d_I
\end{equation}
with
\begin{equation}\label{eq:8}
\Psi_I(x)=p^{[p]T}(x_I-x)c(x)\varphi(x_I-x)
\end{equation}
where $d_I$ is the nodal displacement coefficient vector related to the node $x_I$. $\Psi_I(x)$ is the RK shape function associated with the node $x_I$ and the compact support domain $\supp(x_I)$, $\Omega\in\bigcup_{I=1}^{NP}\supp(x_I)$.
And $p^{[p]}(x)$ is the $p$-degree monomial basis vector.
It is noted that the compact support $\supp(x_I)$ comes from the kernel function $\varphi(x_I-x)$.
Specifically, the cubic and quartic B-spline functions for the construction of quadratic, cubic and quartic meshfree approximants can be stated as follows in 1D case:
\begin{equation}\label{eq:9}
\varphi(r)=\begin{cases}
\dfrac{2}{3}-4r^{2}+4r^{3}, & r\le\dfrac{1}{2}\\[8pt]
\dfrac{4}{3}-4r+4r^{2}-\dfrac{4}{3}r^{3}, & \dfrac{1}{2}<r\le1\\[8pt]
0, & r>1
\end{cases}
\end{equation}
\begin{equation}\label{eq:10}
\varphi(r)=\begin{cases}
\dfrac{115}{192}-\dfrac{625}{64}r^{2}+\dfrac{125}{32}r^{4}, & r\le\dfrac{1}{5}\\[8pt]
\dfrac{55}{96}-\dfrac{625}{96}r^{2}+\dfrac{625}{48}r^{3}-\dfrac{125}{48}r^{4}, & \dfrac{1}{5}<r\le\dfrac{3}{5}\\[8pt]
\dfrac{5}{96}(1-r)^{4}, & \dfrac{3}{5}<r\le1\\[8pt]
0, & r>1
\end{cases}
\end{equation}
with $r=|x_I-x|/s$ being the relative distance between neighboring nodes, $s$ is the size of influence domain.
The 2D kernel function $\varphi(x_I-x)$ can be conveniently defined as the tensor product of two 1D kernel functions in each direction.

The unknown vector $c(x)$ in Eq. \eqref{eq:8} is determined by the following $p$th order reproducing or consistency conditions:
\begin{equation}\label{eq:11}
\sum_{I=1}^{NP}\Psi_I(x)p^{[p]}(x_I-x)=p^{[p]}(0)
\end{equation}
Substituting Eq. \eqref{eq:8} into Eq. \eqref{eq:11}, $c(x)$ takes the following form:
\begin{equation}\label{eq:12}
c(x)=A^{-1}(x)p^{[p]}(0)
\end{equation}
where $A(x)$ is the so-called moment matrix:
\begin{equation}\label{eq:13}
A(x)=\sum_{I=1}^{NP}p^{[p]}(x_I-x)p^{[p]T}(x_I-x)\varphi(x_I-x)
\end{equation}
Then, combining Eq. \eqref{eq:12} with Eq. \eqref{eq:13} yields the following expression of RK shape functions:
\begin{equation}\label{eq:14}
\Psi_I(x)=p^{[p]T}(0)A^{-1}(x)p^{[p]}(x_I-x)\varphi(x_I-x)
\end{equation}

\subsection{Equivalent Galerkin meshfree formulation with Nitsche's method}\label{sec:nitsche}

For convenience of subsequent development, introducing the definition of Sobolev space:
\begin{equation}\label{eq:15}
H_l(\Omega)=\left\{u\mid \mathrm{D}^{|\alpha|}u\in H_0(\Omega),\;|\alpha|\le l\right\}
\end{equation}
\begin{equation}\label{eq:16}
H_0(\Omega)=\left\{u\;\middle|\;\int_{\Omega}u^{2}\dd\Omega<\infty\right\}
\end{equation}
in which $\mathrm{D}^{|\alpha|}(\cdot)$ is the differential operator.

The weak form of Eq. \eqref{eq:6} can be further expressed as following variational form of Kirchhoff--Love plate:
\begin{equation}\label{eq:17}
\text{Find }w\in H_2(\Omega):\quad \A(v,w)=\F(v),\quad \forall v\in H_2(\Omega)
\end{equation}
in which $v$ denotes the test functions, $\A:H_2\times H_2\to\mathbb{R}$ is a bilinear form and $\F:H_2\to\mathbb{R}$ is a linear form.
Assuming that all boundary conditions are essential types ($\Gamma^{w}=\Gamma^{\theta}=\partial\Omega$, $\Gamma^{V}=\Gamma^{M}=\varnothing$), $\A$ and $\F$ in Eq. \eqref{eq:17} can be stated as follows:
\begin{equation}\label{eq:18}
\A(v,w)=\bar{\A}(v,w)+\A_{\Gamma}(v,w),\quad \F(v)=\bar{\F}(v)+\F_{\Gamma}(v)
\end{equation}
where $\A_{\Gamma}(v,w)$ and $\F_{\Gamma}(v)$ are extra terms for imposing essential boundary conditions.
$\bar{\A}$ and $\bar{\F}$ are given as:
\begin{equation}\label{eq:19}
\bar{\A}(v,w)=\int_{\Omega}v_{,ij}C_{ijkl}w_{,kl}\dd\Omega,\quad
\bar{\F}(v)=\int_{\Omega}vq\dd\Omega+\left[\delta w\bar{F}_{c}\right]_{\Gamma}
\end{equation}
Following Cauchy--Schwarz inequality and Nitsche's method, the continuity of $\A(v,w)$ and $\F(v)$ is ensured:
\begin{equation}\label{eq:20}
\left|\A(v,w)\right|\le\gamma\|v\|_{H_2}\|w\|_{H_2},\quad \forall v,w\in H_2(\Omega)
\end{equation}
\begin{equation}\label{eq:21}
\left|\F(v)\right|\le|v|_{H_1}\le\gamma\|v\|_{H_2},\quad \forall v\in H_2(\Omega)
\end{equation}
Also, $\A(v,v)$ is coercive:
\begin{equation}\label{eq:22}
\A(v,v)\ge\alpha\|v\|^{2}_{H_2},\quad \forall v\in H_2(\Omega)
\end{equation}
in which $\alpha$ and $\gamma$ both are positive constants.

In order to imposing the essential boundary conditions, following Nitsche's method, the bilinear form $\A_{\Gamma}(v,w)$ and linear form $\F_{\Gamma}(v)$ in Eq. \eqref{eq:18} can be given as follows:
\begin{equation}\label{eq:23}
\begin{aligned}
\A_{\Gamma}(v,w)&=\int_{\Gamma}V(v)w\dd\Gamma+\int_{\Gamma}vV(w)\dd\Gamma+\beta_1\int_{\Gamma}vw\dd\Gamma\\
&\quad-\int_{\Gamma}M_n(v)w_{,n}\dd\Gamma-\int_{\Gamma}v_{,n}M_n(w)\dd\Gamma+\beta_2\int_{\Gamma}v_{,n}w_{,n}\dd\Gamma
\end{aligned}
\end{equation}
\begin{equation}\label{eq:24}
\begin{aligned}
\F_{\Gamma}(v)&=\int_{\Gamma}V(v)\bar{w}\dd\Gamma+\beta_1\int_{\Gamma}v\bar{w}\dd\Gamma\\
&\quad-\int_{\Gamma}M_n(v)\bar{w}_{,n}\dd\Gamma+\beta_2\int_{\Gamma}v_{,n}\bar{w}_{,n}\dd\Gamma
\end{aligned}
\end{equation}
where $\beta_1$ and $\beta_2$ are stabilized constants, it is noted that the continuity and coercivity conditions still can be ensured within stabilized constants.
$V(v)$ and $M_n(v)$ can be defined as:
\begin{equation}\label{eq:25}
\begin{cases}
q=-M_{ij,ij}=C_{ijkl}w_{,klij}\\[2pt]
V=M_{ij,j}n_i+M_{ns,s}=-C_{ijkl}w_{,kl}n_j-C_{ijkl}w_{,klm}s_ms_j\\[2pt]
M_n=M_{ij}n_in_j=-C_{ijkl}w_{,kl}n_in_j
\end{cases}
\end{equation}

\subsection{Discrete Galerkin meshfree formulation}\label{sec:discrete}

Substituting meshfree approximation into Eq. \eqref{eq:17} generates the discrete meshfree formulation for Kirchhoff--Love plate:
\begin{equation}\label{eq:26}
\text{Find }w^{h}\in V^{h}:\quad \A(v^{h},w^{h})=\F(w^{h}),\quad \forall v^{h}\in V^{h}
\end{equation}
where $V^{h}=\sk\{\Psi_I(x),\,I=1,\ldots,NP\}$ represents the meshfree space, $w^{h}$ is the numerical solution of Eq. \eqref{eq:17}, assuming that the numerical integration is completely accurate.
$\A(v^{h},w^{h})$ and $\F(w^{h})$ are given by:
\begin{equation}\label{eq:27}
\A(v^{h},w^{h})=\sum_{I,J=1}^{NP}v_I\left[K_{IJ}^{[1]}+K_{IJ}^{n_1}+K_{IJ}^{\beta_1}+K_{IJ}^{n_2}+K_{IJ}^{\beta_2}\right]w_J
\end{equation}
\begin{equation}\label{eq:28}
\F(v^{h})=\sum_{I=1}^{NP}v_I\left[f_I+f_I^{n_1}+f_I^{\beta_1}+f_I^{n_2}+f_I^{\beta_2}\right]
\end{equation}
where

(1) $K_{IJ}$, $f_I$ in $\Omega$:
\begin{equation}\label{eq:29}
K_{IJ}=\int_{\Omega}B_I^{T}DB_J\dd\Omega,\quad f_I=\int_{\Omega}\Psi_Iq\dd\Omega
\end{equation}

(2) $K_{IJ}^{n_1}$, $K_{IJ}^{\beta_1}$, $f_I^{n_1}$ and $f_I^{\beta_1}$ on $\Gamma^{u}$:
\begin{equation}\label{eq:30}
\begin{aligned}
K_{IJ}^{n_1}&=-\int_{\Gamma^{u}}\Psi_I\bar{D}C_J\dd\Gamma-\int_{\Gamma^{u}}C_I^{T}\bar{D}^{T}\Psi_J\dd\Gamma,\quad
\bar{D}=h^{T}\left(s_xD_S^{x}+s_yD_S^{y}\right)+n^{T}D_V\\
f_I^{n_1}&=-\int_{\Gamma^{u}}C_I^{T}\bar{D}^{T}\bar{w}\dd\Gamma,\quad
f_I^{\beta_1}=\beta_1\int_{\Gamma^{w}}\Psi_I\bar{w}\dd\Gamma,\quad
K_{IJ}^{\beta_1}=\beta_1\int_{\Gamma^{w}}\Psi_I\Psi_J\dd\Gamma
\end{aligned}
\end{equation}

(3) $K_{IJ}^{n_2}$, $K_{IJ}^{\beta_2}$, $f_I^{n_2}$ and $f_I^{\beta_2}$ on $\Gamma^{\theta}$:
\begin{equation}\label{eq:31}
\begin{aligned}
K_{IJ}^{n_2}&=-\int_{\Gamma^{\theta}}n^{T}A_INDB_J\dd\Gamma-\int_{\Gamma^{\theta}}B_I^{T}DN^{T}A_J^{T}n\dd\Gamma\\
f_I^{n_2}&=-\int_{\Gamma^{\theta}}B_I^{T}DN^{T}\bar{w}_{,n}\dd\Gamma,\quad
f_I^{\beta_2}=\beta_2\int_{\Gamma^{\theta}}n^{T}A_I\bar{w}_{,n}\dd\Gamma,\quad
K_{IJ}^{\beta_2}=\beta_2\int_{\Gamma^{\theta}}n^{T}A_IA_J^{T}n\dd\Gamma
\end{aligned}
\end{equation}
with
\begin{equation}\label{eq:32}
h=\begin{bmatrix}s_xn_x\\ s_yn_y\\ s_xn_y+s_yn_x\end{bmatrix},\quad
B_I=\begin{bmatrix}\Psi_{I,xx}\\ \Psi_{I,yy}\\ 2\Psi_{I,xy}\end{bmatrix},\quad
C_I=\begin{bmatrix}\Psi_{I,xxx}\\ \Psi_{I,xxy}\\ \Psi_{I,xyy}\\ \Psi_{I,yyy}\end{bmatrix},\quad
N=\begin{bmatrix}n_x\\ n_y\\ n_x\\ n_y\end{bmatrix},\quad
\mathbb{N}=\begin{bmatrix}n_x^{2}\\ n_y^{2}\\ 2n_xn_y\end{bmatrix}
\end{equation}
\begin{equation}\label{eq:33}
A=\begin{bmatrix}\Psi_{I,x}\\ \Psi_{I,y}\end{bmatrix},\quad
n=\begin{bmatrix}n_x\\ n_y\end{bmatrix},\quad
D=\begin{bmatrix}1&\nu&0\\ \nu&1&0\\ 0&0&(1-\nu)/2\end{bmatrix}
\end{equation}
and
\begin{equation}\label{eq:34}
D_{\nu}=\begin{bmatrix}1&0&\nu&0\\ 0&\nu&0&1\\ 0&0&1-\nu&0\\ 0&1-\nu&0&0\end{bmatrix},\quad
D_S^{x}=\begin{bmatrix}1&0&\nu&0\\ \nu&0&1&0\\ 0&1-\nu&0&0\end{bmatrix},\quad
D_S^{y}=\begin{bmatrix}0&1&0&\nu\\ 0&\nu&0&1\\ 0&0&1-\nu&0\end{bmatrix}
\end{equation}
Accounting for the effect of numerical integration, Eq. \eqref{eq:26} can be recast as:
\begin{equation}\label{eq:35}
\text{Find }\bar{w}^{h}\in V^{h}:\quad \bA^{*}(v^{h},\bar{w}^{h})=\bF^{*}(v^{h}),\quad \forall v^{h}\in V^{h}
\end{equation}
where $\bar{w}^{h}$ is the final solution of numerically discrete Galerkin meshfree formulation for Kirchhoff--Love plate, $\bA^{*}$ and $\bF^{*}$ respectively represent the discrete $\A$ and $\F$ stated in Eqs. \eqref{eq:27} and \eqref{eq:28}.

\subsection{Reproducing kernel gradient smoothing for Kirchhoff--Love plate}\label{sec:rkgs}

Assuming the moment and loading can be spanned by the basis vector:
\begin{equation}\label{eq:36}
M_{ij}=p^{[p-2]T}(x)\alpha_{ij},\quad q=-p^{[p-2]T}_{,ij}(x)\alpha_{ij}
\end{equation}
Introducing Eq. \eqref{eq:36} into Eq. \eqref{eq:6} obtains:
\begin{equation}\label{eq:37}
\begin{aligned}
\int_{\Omega}\delta w_{,ij}p^{[p-2]T}\dd\Omega
={}&-\int_{\Gamma}\delta w\left[\begin{matrix}p^{[p-2]T}_{,s}\left(n_is_j+n_js_i\right)/2\\[2pt]
+\left(p^{[p-2]T}_{,j}n_i+p^{[p-2]T}_{,i}n_j\right)/2\end{matrix}\right]\dd\Gamma
+\int_{\Omega}\delta wp^{[p-2]T}_{,ij}\dd\Omega\\
&+\int_{\Gamma}\delta w_{,n}p^{[p-2]T}n_in_j\dd\Gamma
+\left[\delta wp^{[p-2]T}\left(n_is_j+n_js_i\right)/2\right]_{\Gamma}
\end{aligned}
\end{equation}
Further introducing meshfree approximation into Eq. \eqref{eq:37}, then we attain the integration constraint for Kirchhoff--Love plate:
\begin{equation}\label{eq:38}
\begin{aligned}
\int_{\Omega}\Psi_{I,ij}p^{[p-2]T}\dd\Omega
={}&-\int_{\Gamma}\Psi_I\left[\begin{matrix}p^{[p-2]T}_{,s}\left(n_is_j+n_js_i\right)/2\\[2pt]
+\left(p^{[p-2]T}_{,j}n_i+p^{[p-2]T}_{,i}n_j\right)/2\end{matrix}\right]\dd\Gamma
+\int_{\Omega}\Psi_Ip^{[p-2]T}_{,ij}\dd\Omega\\
&+\int_{\Gamma}\Psi_{I,m}n_mp^{[p-2]T}n_in_j\dd\Gamma
+\left[\delta wp^{[p-2]T}\left(n_is_j+n_js_i\right)/2\right]_{\Gamma}
\end{aligned}
\end{equation}
In accordance with reproducing kernel gradient smoothing framework \cite{ref43}, the second order smoothed gradient $\tilde{\Psi}_{I,ij}$ can be expressed as follows:
\begin{equation}\label{eq:39}
\tilde{\Psi}_{I,ij}=p^{[p-2]T}(x)c_{ij}^{[p-2]}(x_I)\hat{\varphi}(x)
\end{equation}
where $c_{ij}^{[p-2]}(x_I)$ is an unknown vector to be determined.
$\hat{\varphi}(x)$ is a kernel function that can be conveniently selected as the following simple form:
\begin{equation}\label{eq:40}
\hat{\varphi}(x)=\begin{cases}
1, & x\in\Omega_C\\[2pt]
0, & x\notin\Omega_C
\end{cases}
\end{equation}
where $\Omega_C$ is a generic non-overlapping integration cell and we have $\bigcup_{C=1}^{NC}\Omega_C=\Omega$, $NC$ denotes the total number of integration cells.

Substituting Eq. \eqref{eq:39} into Eq. \eqref{eq:38} within a integration cell level gives:
\begin{equation}\label{eq:41}
G_Cc_{ij}^{[p-1]}(x_I)=g_{ijI}^{C}
\end{equation}
with
\begin{equation}\label{eq:42}
G_C=\int_{\Omega_C}p^{[p-2]}p^{[p-2]T}\dd\Omega
\end{equation}
\begin{equation}\label{eq:43}
\begin{aligned}
g_{ijI}^{C}={}&-\int_{\Gamma_C}\Psi_I\left[\begin{matrix}p^{[p-2]T}_{,s}\left(n_is_j+n_js_i\right)/2\\[2pt]
+\left(p^{[p-2]T}_{,j}n_i+p^{[p-2]T}_{,i}n_j\right)/2\end{matrix}\right]\dd\Gamma
+\int_{\Omega_C}\Psi_Ip^{[p-2]T}_{,ij}\dd\Omega\\
&+\int_{\Gamma_C}\Psi_{I,m}n_mp^{[p-2]T}n_in_j\dd\Gamma
+\sum_{m=1}^{n_c}\left\{\begin{matrix}\left[\Psi_Ip^{[p-2]T}\left(n_is_j+n_js_i\right)/2\right](c_{m+1})\\[4pt]
-\left[\Psi_Ip^{[p-2]T}\left(n_is_j+n_js_i\right)/2\right](c_m)\end{matrix}\right\}
\end{aligned}
\end{equation}
It is noted that the gradient smoothing moment matrix $G_C$ of Eq. \eqref{eq:42} is essentially a constant matrix within each integration cell $\Omega_C$.

From Eq. \eqref{eq:41}, we immediately have $c_{ij}^{[p-1]}(x_I)=G_C^{-1}g_{ijI}^{C}$ and accordingly the second order smoothed gradients of meshfree shape functions become:
\begin{equation}\label{eq:44}
\tilde{\Psi}_{I,ij}=p^{[p-1]T}(x)G_C^{-1}g_{ijI}^{C}
\end{equation}
By construction, the proposed reproducing kernel smoothed gradient $\tilde{\Psi}_{I,ij}$ given by Eqs. \eqref{eq:44} identically meet the integration constraint of Eq. \eqref{eq:38}.
This fact can be straightforwardly verified as follows:
\begin{equation}\label{eq:45}
\begin{aligned}
\int_{\Omega}\tilde{\Psi}_{I,ij}p^{[p-2]T}\dd\Omega
&=\sum_{C=1}^{NC}\int_{\Omega_C}p^{[p-2]}p^{[p-2]T}G_C^{-1}g_{ijI}^{C}\dd\Omega\\
&=\sum_{C=1}^{NC}\underbrace{\left[\int_{\Omega_C}p^{[p-2]}p^{[p-2]T}\dd\Omega\right]G_C^{-1}g_{ijI}^{C}}_{=\,G_C}
=\sum_{C=1}^{NC}g_{ijI}^{C}=g_{ijI}
\end{aligned}
\end{equation}
